\documentclass[10pt,a4paper]{amsart}
\usepackage[margin=19mm]{geometry}
\usepackage{amsmath,amssymb,mathtools}
\usepackage[scr=rsfs,cal=euler]{mathalfa}
\usepackage{xfrac}
\usepackage[unicode,hidelinks,bookmarks=false]{hyperref}
\newcommand{\ie}{i.e.\ }
\newcommand{\cf}{cf.\ }
\newcommand{\resp}{respectively\ }
\newcommand{\Subsection}{\subsection}
\providecommand{\nobreakdash}{\nobreak\hbox{-}}
\setlength{\emergencystretch}{2em}
\multlinegap0pt
\usepackage{bm}
\usepackage{bbm}
\usepackage{dsfont}
\usepackage{xspace}
\usepackage{cancel}
\usepackage{mathtools}
\usepackage{amsthm}
\makeatletter
\@ifundefined{theorem}{\newtheorem{theorem}{Theorem}}{}
\@ifundefined{definition}{\newtheorem{definition}[theorem]{Definition}}{}
\makeatother
\newcommand{\COL}{\operatorname{COL}}
\newcommand{\Commonconst}{A}
\newcommand{\N}{\mathbb{N}}
\newcommand{\commonconst}{a}
\newcommand{\numcolors}{c}
\newcommand{\numpoints}{n}
\newcommand{\oequiv}{\approx}
\newcommand{\om}{\omega}
\newcommand{\rareconst}{b}
\newcommand{\setconcat}{+}
\title[Big Ramsey Degrees of Countable Ordinals]{Big Ramsey Degrees of Countable Ordinals}
\author{Joanna Boyland, William Gasarch, Nathan Hurtig, Robert Rust}
\date{Source: arXiv:2305.07192v3 (CC BY 4.0)}
\begin{document}
\maketitle

\noindent\emph{Source and licence.} Big Ramsey Degrees of Countable Ordinals. Version arXiv:2305.07192v3 (CC BY 4.0), distributed under the Creative Commons Attribution 4.0 International licence, as recorded in the arXiv licence metadata for this exact version and shown on the abstract page https://arxiv.org/abs/2305.07192v3. Full rights notice: licenses/arxiv-2305.07192-CC-BY-4.0.txt.\par\smallskip

\noindent\textit{Editorial scope.} Counted unit: the Theorem whose source label is \texttt{th:om2} in the article (source0.tex lines 706--708, source label \texttt{th:om2}), delivered together with the article's own complete original proof of it (source0.tex lines 709--803, one \texttt{proof} environment of 95 lines). No mathematics is written, repaired or re-derived: statement and proof are the article's own text line for line. Required context retained uncounted: th:ramsey (lines 241--251); def:coef (lines 699--702). Every definition, display and label of those lines is retained unchanged. Omitted from this package and not redistributed: 1--240, 252--698, 703--705, 804--1919. Nothing inside a retained excerpt is omitted or altered: every retained body is delivered exactly as the article's lines are written, so no omission receipt of any kind accompanies this package. Citations: the retained excerpts carry no citation command at all, so no bibliography entry of the article is transcribed and no reference is redistributed. Reference adaptation: none was needed and none was made; every reference inside the delivered unit resolves inside it, so no unresolved reference or citation is produced. Printed slips kept literal: no printed word, symbol, hypothesis or number of the counted statement or of its proof was altered. Layout and class: the article's own class is \texttt{article} with packages including amsthm, amssymb, multirow, amsmath, hyperref, geometry, inputenc; the wrapper is the lane's pinned \texttt{amsart} editorial wrapper at 10pt on A4, which restates the theorem-like environments the retained text needs and adds the article's own macro definitions that the retained text needs (\texttt{\textbackslash COL}, \texttt{\textbackslash Commonconst}, \texttt{\textbackslash N}, \texttt{\textbackslash commonconst}, \texttt{\textbackslash numcolors}, \texttt{\textbackslash numpoints}, \texttt{\textbackslash oequiv}, \texttt{\textbackslash om}, \texttt{\textbackslash rareconst}, \texttt{\textbackslash setconcat}), with extra wrapper packages (bm, bbm, dsfont, xspace, cancel, mathtools). The delivered numbering is the unit's own. Assets: no retained span contains a figure environment, a graphics-inclusion command, a PDF-page inclusion, a table or a TikZ picture; no image file is copied into this package, the package declares no asset dependency and no image file is redistributed with it. Licence: version arXiv:2305.07192v3 of this article is distributed under the Creative Commons Attribution 4.0 International licence, as recorded in the arXiv licence metadata for this exact version and shown on the abstract page https://arxiv.org/abs/2305.07192v3. A full rights notice is delivered at licenses/arxiv-2305.07192-CC-BY-4.0.txt. Characters: every retained excerpt is pure ASCII, so the delivered engine, XeLaTeX with its default Latin Modern text font, reports no missing character.
\begin{theorem}\label{th:ramsey}
For all $\numpoints \in \N,$
% \begin{enumerate}
    % \item 
$T(\numpoints,\omega)=1.$
% \item 
% Let $\numcolors \geq 1$ and $\COL\colon\binom{\omega}{\numpoints}\to [\numcolors]$.
% Then there exists some $H\oequiv \omega$ such that 
% $$\left | \COL\left(\binom{H}{\numpoints}\right) \right | = 1.$$
% \end{enumerate}
\end{theorem}

\begin{definition} \label{def:coef}
If $e \in \binom{\om^d}{n},$ we refer to $e$ as an {\it edge.} Then $e$ has $n$ {\it elements} selected from $\om^d.$
Hence each element of $e$ is some ordinal number less than $\om^d.$ That element can be represented as a polynomial of degree $< d$ in $\omega.$ We refer to the natural-valued coefficients of these polynomials as the {\it coefficients} of $e.$
\end{definition}

\begin{theorem}\label{th:om2} 
The big Ramsey degree $T(2,\om^2)=4$.
\end{theorem} 

\begin{proof}
We first prove $T\left(2, \om^2\right) \leq 4$.

Let $\numcolors \in \N$ and let \begin{align*}
    \COL \colon \binom{\om^2}{2} \to [\numcolors]
\end{align*}
be a $\numcolors$-coloring of $\binom{\om^2}{2}$. 
We define four functions $f_0,f_1,f_2,f_3$ from 
domain $\binom{\om}{4}$ to codomain $\binom{\om^2}{2}$ and then use them to define a coloring 
from 
$\binom{\om}{4}$ to $[\numcolors]^4$.
In what follows, 
let $x_1<x_2<x_3<x_4$. For $i \in [4],$ we define $f_i \colon \binom{\om}{4} \to \binom{\om^2}{2}$ where \begin{align*}
f_0(x_1,x_2,x_3,x_4)  &= \{\om \cdot x_1 + x_2, \om \cdot x_3 + x_4\}\\
f_1(x_1,x_2,x_3,x_4)  &= \{\om \cdot x_1 + x_3, \om \cdot x_2 + x_4\}\\
f_2(x_1,x_2,x_3,x_4)  &= \{\om \cdot x_1 + x_4, \om \cdot x_2 + x_3\}\\
f_3(x_1,x_2,x_3,x_4)  &= \{\om \cdot x_1 + x_2, \om \cdot x_1 + x_3\}.
\end{align*}

Then let
$\COL' \colon \binom{\om}{4} \to [\numcolors]^4$ where
$$\COL'(X) =(\COL(f_0(X)),
\COL(f_1(X)),
\COL(f_2(X)),
\COL(f_3(X))
).
$$

Apply Theorem~\ref{th:ramsey} on $\COL'$ to find some $G \oequiv \omega$ where 
$\left|\COL' \left( \binom{G}{4} \right)\right|=1$.
Let $G = \{\commonconst_0 <\nobreak \commonconst_1 <\nobreak \cdots\}.$ Let
\begin{align*}
H = & ~~~~\!~\{\om \cdot \commonconst_{1} + \commonconst_2, \om \cdot \commonconst_{1} + \commonconst_6,~\om \cdot \commonconst_{1} + \commonconst_{10}, \ldots\} \\
 & \setconcat ~\{\om \cdot \commonconst_{3} + \commonconst_{4}, \om \cdot \commonconst_{3} + \commonconst_{12}, \om \cdot \commonconst_{3} + \commonconst_{20}, \ldots\}\\
 & \setconcat~\{\om \cdot \commonconst_{5} + \commonconst_{8}, \om \cdot \commonconst_{5} + \commonconst_{24}, \om \cdot \commonconst_{5} + \commonconst_{40}, \ldots\}\\
& ~~\!\vdots
\end{align*}
Formally, $$
H = \Commonconst_1 \setconcat \Commonconst_2 \setconcat \cdots
$$ where $$
\Commonconst_i = \{\om\cdot \commonconst_{2i-1} + \commonconst_j \colon j = 2^i + k 2^{i+1}, k \in \N\}.
$$
Note that each of the $A_i$'s coefficients
(as defined in Definition~\ref{def:coef}) 
are disjoint.

Then $H \oequiv \om^2$, as it is the concatenation of countably infinite sets order-equivalent to $\om$. Note that for every $\om \cdot \commonconst_i + \commonconst_j \in H$ we have $\commonconst_i < \commonconst_j$.

Let the single 4-tuple that $\COL'$ outputs on $\binom{G}{4}$ be $Y.$ Consider any edge $\{\om \cdot \rareconst_1 + \rareconst_2, \om \cdot \rareconst_3 + \rareconst_4\} \in \binom{H}{2}$ with $\om \cdot \rareconst_1 + \rareconst_2 < \om \cdot \rareconst_3 + \rareconst_4$. \begin{itemize}
\item If $\rareconst_1 \neq \rareconst_3$, then the two elements are from different $A_i$ sets. Then $\rareconst_1 < \rareconst_3$ by the ordering of the two elements and $\rareconst_2 \neq \rareconst_4$ by the construction of $H$. 
We also have $\rareconst_1 < \rareconst_2$, $\rareconst_1 < \rareconst_4$, and $\rareconst_3 < \rareconst_4$ by the construction of $H$.
There are three subcases: \begin{itemize}
    \item $\rareconst_1<\rareconst_2<\rareconst_3<\rareconst_4$: then $\COL(f_0(\rareconst_1,\rareconst_2,\rareconst_3,\rareconst_4)) \in Y.$
    \item $\rareconst_1<\rareconst_3<\rareconst_2<\rareconst_4$: then $\COL(f_1(\rareconst_1,\rareconst_3,\rareconst_2,\rareconst_4)) \in Y.$
    \item $\rareconst_1<\rareconst_3<\rareconst_4<\rareconst_2$: then $\COL(f_2(\rareconst_1,\rareconst_3,\rareconst_4,\rareconst_2)) \in Y.$
\end{itemize}
In all subcases, $\COL(\{\om \cdot \rareconst_1 + \rareconst_2, \om \cdot \rareconst_3 + \rareconst_4\}) \in Y.$
% Then either $\rareconst_1<\rareconst_2<\rareconst_3<\rareconst_4$, $\rareconst_1<\rareconst_3<\rareconst_2<\rareconst_4$, or $\rareconst_1<\rareconst_3<\rareconst_4<\rareconst_2$. In each of the three cases, $f_1(\rareconst_1,\rareconst_2,\rareconst_3,\rareconst_4) \in Y$, $f_2(\rareconst_1,\rareconst_3,\rareconst_2,\rareconst_4) \in Y$, and $f_3(\rareconst_1,\rareconst_3,\rareconst_4,\rareconst_2) \in Y$ respectively so $\COL(\{\om \cdot \rareconst_1 + \rareconst_2, \om \cdot \rareconst_3 + \rareconst_4\}) \in Y$.
\item If $\rareconst_1 = \rareconst_3$, then the two elements are from the same $A_i.$ Then $\rareconst_2 < \rareconst_4$ by the ordering of the elements and so $\rareconst_1=\rareconst_3<\rareconst_2<\rareconst_4$ by the construction of $H$. Because $\COL(f_3(\rareconst_1,\rareconst_2,\rareconst_4,\rareconst_4+1)) \in Y$ (note that the output of $f_3$ does not depend on its final argument), $\COL(\{\om \cdot \rareconst_1 + \rareconst_2, \om \cdot \rareconst_3 + \rareconst_4\}) \in Y$.
\end{itemize}
In all cases, $\COL(\{\om \cdot \rareconst_1 + \rareconst_2, \om \cdot \rareconst_3 + \rareconst_4\}) \in Y$ so \begin{align*}
\COL\left(\binom{H}{2}\right) \subseteq Y
\end{align*} with $|Y|=4$. Because $H \oequiv \om^2$ and $\COL$ was arbitrary, $T\left(2, \om^2\right) \leq 4$.

We now prove $T\left(2, \om^2\right) \geq  4$.
Let $\COL \colon \binom{\om^2}{2} \to [4]$ with \begin{align*}
\COL(\om \cdot \commonconst_1 + \commonconst_2, \om \cdot \commonconst_3 + \commonconst_4) = 
\begin{cases}
0 & \commonconst_1 < \commonconst_2 < \commonconst_3 < \commonconst_4\\
1 & \commonconst_1 < \commonconst_3 < \commonconst_2 < \commonconst_4\\
2 & \commonconst_1 < \commonconst_3 < \commonconst_4 < \commonconst_2\\
3 & \textnormal{otherwise}
\end{cases}
\end{align*}
where $\om \cdot \commonconst_1 + \commonconst_2 < \om \cdot \commonconst_3 + \commonconst_4$. 
Let $H \subseteq \omega^2$ be any set such that $H\oequiv \omega^2$. 
We show that $\left|\COL\left(\binom{H}{2}\right)\right)|=4$. Every element of $H$ is of the form $\omega \cdot \rareconst + e,$ for finite $\rareconst,e.$ Because $H \oequiv \om^2$, we have $H = H_1+H_2+ \cdots$ where each $H_i \oequiv \omega$. For each $H_i,$ there exists some $\rareconst_i$ such that $$
H'_i = \{\omega \cdot \rareconst_i + e \in H_i\} \oequiv \om.
$$ Let $e_{i,j}$ be such that every element of $H'_i$ is of the form $\omega \cdot \rareconst_i + e_{i,j},$
where
$$\rareconst_1 < \rareconst_2 < \cdots$$
and for all $i$, 
$$e_{i,1} < e_{i,2} < \cdots.$$
We now show each color is output by $\COL$ on $H$.
\begin{itemize}
    \item Color 0: Starting with any $\rareconst_i$, find some $j$ where $e_{i, j} > \rareconst_i$. This is guaranteed, as $\rareconst_i$ is finite and the $e_{i, j}$ are infinitely increasing. Then find some $k$ where $\rareconst_k > e_{i, j}$; again guaranteed because $e_{i, j}$ is finite and the $\rareconst_i$ are infinitely increasing. Finally, find an
    $\ell$ where $e_{k ,\ell} > \rareconst_k$ guaranteed by similar means. Then 
    $$\COL(\omega\cdot \rareconst_i + e_{i, j}, \omega\cdot \rareconst_k + e_{k, \ell})=0.$$
    \item Colors 1 and 2 are guaranteed to exist by arguments similar to the argument for color 0.
    \item Color 3: The case of $\commonconst_1=\commonconst_3<\commonconst_2<\commonconst_4$ falls into the category of {\it otherwise}. With this in mind, we can start with some $\rareconst_i$, and then find a $e_{i, j} > \rareconst_i$. Then, we only need to find a $e_{i, \ell} > e_{i, j}$; this is guaranteed because $e_{i, j}$ is finite. Then $$\COL(\omega\cdot \rareconst_i + e_{i, j}, \omega\cdot \rareconst_i + e_{i, \ell})=3.$$
\end{itemize}

% While it was simple to find edges that output each of these four colors, the proof that $T(2,\om^2) \leq 4$ shows that a coloring with more than four colors can never guarantee that more than four colors are output on an order-equivalent subset.
 
\end{proof}

\end{document}
